Numerous studies have delved into enhancing energy efficiency within telecommunication networks. This research can be categorized into two main areas: methods for measuring energy consumption and strategies for developing algorithms and methodologies to optimize it. However, to date, there is a gap in understanding the impact of realistic policies in real-world settings. Realistically assessing how different policies perform requires parallel comparisons, akin to A/B testing in the web domain, to evaluate the energy gains or savings of one policy against another.

\textbf{Measuring energy consumption.} The energy consumption of mobile networks is one of the major concerns of the telecom industry. Piovesan et al. \cite{Piovesan2022} proposes a model for characterization of the power consumption of 5G multi-carrier BSs, build on a large data collection campaign. Additionally, Lopes et al. \cite{survey_5g_energy_efficiency} conducted a survey outlining the primary energy efficiency technologies provided by 3GPP NR, together with their main benefits and challenges. Moreover, the community acknowledges the need for sustainability, with efforts underway to quantify the carbon footprint of communication networks \cite{IEA2024}. In the computer systems community, there are existing approaches to quantify and optimize carbon emissions, \cite{CarbonScaler}, by exploit the temporal flexibility inherent to many cloud workloads. In this study, we present utilization and performance data gathered at various times throughout the day. This information can be instrumental in assessing whether base station utilization can indeed benefit from temporal flexibility, a crucial aspect in achieving greater carbon efficiency.

\textbf{Optimizing energy saving policies.} Energy-saving policies involve the dynamic shutdown of base stations during low-traffic periods and their subsequent reactivation as traffic increases.Accurately predicting traffic loads in the near future is crucial for determining the appropriate times to switch base stations on and off. Donevski et al \cite{Donevski2019} propose using dense Neural Network and Recurrent Neural Network paradigms, although their approach lacks empirical validation through real-world trials. Moreover, Domenico et al. \cite{user_transfer_5g} focus on the carrier shutdown approach, enabling base stations to autonomously power down during low-traffic periods by transferring their load to neighboring active base stations. Lastly, it is crucial to acknowledge that such optimization involves the sudden removal and addition of carriers to the network. Therefore, there is a need to develop automated and precise methods for generating configuration parameters for newly added carriers in cellular networks, as suggested by \cite{Mahimkar2021}, such work is orthogonal to the one we are presenting.